% Esqueleto de memoria de TFG con la clase tfgtfmthesisuam (EPS-UAM).
% Compilar con pdfLaTeX o XeLaTeX (en Overleaf, el compilador por defecto vale).
% Las marcas invisibles para la comprobación automática de TFGTool las pone la clase.
\documentclass[epsbased,copyright,final,printable,extendedindex,firstnumbered,tfg,overleaf]{tfgtfmthesisuam}

\advisor{Nombre del tutor o tutora}
\levelin{Ingeniería Informática}
\title{Título del Trabajo Fin de Grado}
\author{Nombre del autor o autora}
\copyrightdate{Junio de 2027}

\resumenfile{inicio/resumen}
\abstractfile{inicio/abstract}
\keywords{Palabra clave 1, palabra clave 2}
\palabrasclave{Palabra clave 1, palabra clave 2}

\coverdata
{
  Escuela Politécnica Superior \\
  Universidad Autónoma de Madrid \\
  C\textbackslash Francisco Tomás y Valiente nº 11
}

\bibliographyconfig{tfgtfmthesisuam}

\graphicsdir{img}
\logosdir{img}

\begin{document}

% El cuerpo de la memoria (máximo 40 páginas) va desde aquí hasta las conclusiones
\chapter{Introducción\label{CAP:INTRO}}{capitulos/introduccion}
\chapter{Estado del arte\label{CAP:ESTADO}}{capitulos/estado}
\chapter{Diseño y desarrollo\label{CAP:DESARROLLO}}{capitulos/desarrollo}
\chapter{Pruebas y resultados\label{CAP:PRUEBAS}}{capitulos/pruebas}
\chapter{Conclusiones y trabajo futuro\label{CAP:CONCLUSIONES}}{capitulos/conclusiones}

% La bibliografía (y el glosario, si lo hay) se incluyen automáticamente aquí
\appendix

\chapter{Manual de usuario\label{AP:MANUAL}}{apendices/manual}

\end{document}
